\section{Related Work}

\paragraph{Faithfulness of chain-of-thought reasoning.}
Chain-of-thought (CoT) prompting was introduced as a way to elicit
intermediate reasoning steps that improve downstream performance
\citep{wei2022chain}. Subsequent work questioned whether such traces
faithfully represent the model's actual decision process.
\citet{turpin2023language} show that biasing prompt features such as
re-ordering multiple-choice options so that the answer is always (A)
can systematically alter final predictions while leaving the verbalized
CoT largely silent about that influence.
\citet{lanham2023measuring} extend this analysis with causal
interventions on the trace itself, finding that models often produce
plausible-looking reasoning that the final answer does not actually
depend on. \citet{paul2024making} and \citet{arcuschin2025chainofthought}
report similar dissociations across a broader set of tasks, and
\citet{chen2025reasoning} document the same pattern in modern reasoning
models. We adopt the operational definition of faithfulness from this
line of work---whether the model's CoT explicitly references the cause
of a behavioral change---and apply it to multilingual, culturally
grounded inputs.

\paragraph{Sycophancy and biased-hint injection.}
Sycophancy, the tendency of LLMs to align with cues from a user or
prompt even when those cues are wrong, has been documented across
generation, classification, and preference-rating tasks
\citep{perez2023discovering,sharma2023towards,wei2023simple}.
\citet{turpin2023language}'s ``Answer Always (X)'' few-shot biased CoT
is the canonical case study of \emph{behavioral} susceptibility paired
with \emph{verbal} silence. We build directly on this design: we
inject authority, social, indirect, and Turpin-style few-shot biased
cues, then ask whether the model's CoT acknowledges the cue and whether
its final answer flips. Unlike prior work, we vary the language of the
cue independently of the language and culture of the task content.

\paragraph{Culturally grounded NLP and multilingual reasoning.}
A growing body of work shows that LLMs encode culture-specific defaults
that diverge across languages
\citep{masoud2025cultural,naous2024multicultural,hammerl2022moral,aksoy2024morality}.
\citet{shi2024culturebank} introduce CultureBank, a large-scale
community-sourced repository of cultural descriptors which we use to
ground our task items in attested norms rather than stereotype.
On the reasoning side, \citet{shi2023multilingual} demonstrate that CoT
ability transfers across languages, but evaluation has focused on
correctness; the \emph{faithfulness} of multilingual CoT---and its
interaction with culturally targeted bias cues---remains
underexplored. The closest concurrent work,
\citet{fu2025multilingualjudge}, finds that multilingual
LLM-as-a-judge verdicts are unreliable across languages, but that
study evaluates judging behavior rather than the faithfulness of CoT
under hint injection.

\paragraph{Multilingual safety and benchmarks.}
For cross-domain comparison we use GSM8K \citep{cobbe2021gsm8k}
as an objective-answer mathematical reasoning task, and XSAFETY
\citep{wang2024xsafety} as a multilingual safety benchmark.
This lets us ask whether the faithfulness patterns we observe on
culturally grounded items also hold in domains where the
``correct'' answer is either objective (math) or normatively fixed
(refusal). BIG-Bench Hard \citep{suzgun2023bbh,srivastava2022bigbench}
is referenced as a methodological precedent for CoT evaluation.
